problem:  
Let \((\mathcal{L}^{n+1}, \bar{g})\) be a Lorentzian Einstein manifold with \(\operatorname{Ric}_{\bar{g}} = \kappa\, \bar{g}\).  
Let \(M^n \subset \mathcal{L}^{n+1}\) be a complete spacelike hypersurface with induced Riemannian metric \(g\), second fundamental form \(h\), and future-pointing unit normal vector field \(N\).  
Suppose there exists a smooth function \(f : M \to \mathbb{R}\) satisfying the gradient Ricci soliton equation  
\[
\operatorname{Ric}_g + \nabla^2 f = \lambda g,
\]
and that \(f\) admits a smooth extension \(\bar{f}\) to a neighborhood of \(M\) in \(\mathcal{L}^{n+1}\) satisfying:  

1. The ambient gradient field \(\bar{V} = \bar{\nabla}\bar{f}\) is null and tangent to a congruence of null geodesics, i.e. \(\bar{\nabla}_{\bar{V}}\bar{V}=0\) and \(\bar{g}(\bar{V},\bar{V})=0\);  
2. The tangential projection \(V = (\bar{V})^T\) coincides with the intrinsic gradient \(\nabla f\): \(V=\nabla f\).  

**Conjecture / Open Problem:**  
Suppose further that \((\bar{\mathcal{L}},\bar{g},\bar{V})\) admits no parallel null direction other than multiples of \(\bar{V}\).  
Does the existence of such an ambient null-geodesic extension of \(f\) force the projected vector field \(V=\nabla f\) to be an eigenvector of the shape operator \(A\) of \(M\)?  
Equivalently, must the principal curvature directions of \(M\) split such that one eigenspace corresponds to the flow of \(\nabla f\)?  

If so, does this alignment yield (locally) a warped product decomposition  
\[
(M,g) \cong I \times_\phi \Sigma^{n-1},
\]
where the warping function \(\phi\) satisfies an ODE determined by the null geodesic expansion rate of \(\bar{V}\)?  
A positive answer would identify a new rigidity mechanism linking intrinsic Ricci soliton structure to ambient null causality.

Why is it a "good" problem:  
This question is tightly formulated, mathematically coherent, and deals with a setting not covered by existing rigidity results for Ricci solitons. It isolates a single conjectural mechanism—the **alignment of the soliton potential gradient with a principal curvature direction**—that could link intrinsic Riemannian soliton geometry to extrinsic Lorentzian null structure.  
It lies at the intersection of **geometric analysis**, **Lorentzian causality theory**, and **submanifold geometry**.  The hypothesis that an ambient potential has null geodesic gradient introduces causal constraints unknown in the standard Riemannian soliton theory, requiring integration of techniques from both Ricci flow analysis and null hypersurface geometry.  
A resolution would reveal whether lightlike structures in pseudo-Riemannian manifolds impose curvature-alignment rigidity on Ricci solitons, opening a novel connection between Riemannian solitons, Lorentzian geometry, and general relativity. In form, it is simple—the alignment of one vector field with a curvature direction—but in depth, it probes foundational aspects of how curvature, causality, and gradient flows interact.